\documentclass[letterpaper,12pt]{article}
\usepackage[T1]{fontenc}
\usepackage[utf8]{inputenc}
\usepackage[spanish]{babel}
\usepackage{framed}
\usepackage{amsfonts,setspace}
\usepackage{amsmath}
\usepackage{amssymb, amsmath, amsthm}
\usepackage{comment}
\usepackage[most]{tcolorbox}
\usepackage{amssymb}
\usepackage{dsfont}
\usepackage{anysize}
\usepackage{pdfpages}
\usepackage{multicol}
\usepackage{enumerate}
\usepackage{graphicx}
\usepackage{float}
\usepackage[dvipsnames]{xcolor}
\usepackage[left=1.5cm,top=0.9cm,right=1.5cm, bottom=1.4cm]{geometry}
\usepackage{fancyhdr}
\pagestyle{fancy}
\definecolor{blue(pigment)}{rgb}{0.2, 0.2, 0.6}
% Page
\oddsidemargin -0.4in
\textwidth 7in
\topmargin -0.6in
\textheight 9.1in
\parindent 0em
\parskip 1ex

%Teoremas, Lemas, etc.
\theoremstyle{plain}
\newtheorem{teo}{Teorema}
\newtheorem{lem}{Lemma}
\newtheorem{prop}{Proposici\'on}
\newtheorem{cor}{Corolario}
\theoremstyle{definition}
\newtheorem{defi}{Definici\'on}
\newtheorem{eje}{Ejemplo}
\newtheorem{obs}{Observaci\'on}
% fin macros

% Comandos más usados
% Conjuntos
\newcommand{\N}{\mathbb{N}}
\newcommand{\Z}{\mathbb{Z}}
\newcommand{\Q}{\mathbb{Q}}
\newcommand{\R}{\mathbb{R}}
\newcommand{\C}{\mathbb{C}}
\newcommand{\RR}{\mathbb{R} \times \mathbb{R}}
% Flechas
\newcommand{\ssi}{\Longleftrightarrow}
\newcommand{\imp}{\Longrightarrow}
\newcommand{\pmi}{\Longleftarrow}
% Trigonométricas
\newcommand{\sen}{\text{sen}}

\usepackage{versions}
\includeversion{solution}

\let\epsilon\varepsilon

\begin{document}
\def\Pauta{1}
%%%%%ESCRIBIR 1 si quieren ver la Prueba con Pauta%%%%
%%%%%ESCRIBIR 0 si quieren ver sólo la Ayudantía%%%%
	
	\fancyhead[L]{\itshape{Escuela de Ingeniería}}
	\fancyhead[R]{\itshape{Universidad de O'Higgins}}
	
	\begin{minipage}{11.5 cm}
         \vspace{-6mm}
		\begin{flushleft}
			\vspace{1mm}
			\textbf{ING1002-2 Cálculo Diferencial e Integral}\\
			\textbf{Profesor:}  Gonzalo Flores G.  \\
			\textbf{Ayudante:}  Juan I. Riquelme \& Cristian Acevedo R. \\
		\end{flushleft}
	\end{minipage}
	
	\begin{picture}(2,3)
		\put(430,9.5){\includegraphics[scale=0.22]{Imagenes/UOH.png}}
	\end{picture}
	\vspace{-6mm}
	\begin{center}
		\LARGE \bf{Pauta Ayudantía 11 $+$ ejercicios extra} %Teorema Fundamental del Cálculo
	\end{center}
    
\begin{enumerate}[\textbf{P\arabic{enumi}.}]

\item Usando integración por partes calcule las siguientes primitivas
\begin{enumerate}
    \item $\displaystyle \int x^2 \sen(x) \, dx$
    \item $\displaystyle \int x^3 \ln^2(x) dx$
\end{enumerate}

\color{blue}
\textbf{Sol:} a) una vaca sin cola vestida de uniforme
\[
u = x^2 \quad \Rightarrow \quad du = 2x\,dx
\]
\[
dv = \sen(x)\,dx \quad \Rightarrow \quad v = -\cos(x)
\]
\[
\int x^2 \sen(x)\,dx
= -x^2\cos(x) + \int 2x\cos(x)\,dx
\]
\[
\int 2x\cos(x)\,dx
\]
\[
u = 2x \quad \Rightarrow \quad du = 2\,dx
\]
\[
dv = \cos(x)\,dx \quad \Rightarrow \quad v = \sen(x)
\]
\[
\int 2x\cos(x)\,dx
= 2x\sen(x) - \int 2\sen(x)\,dx
\]
\[
= 2x\sen(x) + 2\cos(x)
\]
\[
\int x^2 \sen(x)\,dx
= -x^2\cos(x) + 2x\sen(x) + 2\cos(x) + C
\]
b) para $\displaystyle 
\int x^3 \ln^2(x)\,dx$
\[
u = \ln^2(x) \quad \Rightarrow \quad du = \frac{2\ln(x)}{x}\,dx
\]
\[
dv = x^3\,dx \quad \Rightarrow \quad v = \frac{x^4}{4}
\]
\[
\int x^3 \ln^2(x)\,dx
= \frac{x^4}{4}\ln^2(x) - \int \frac{x^4}{4}\cdot\frac{2\ln(x)}{x}\,dx
\]
\[
= \frac{x^4}{4}\ln^2(x) - \frac{1}{2}\int x^3\ln(x)\,dx
\]
\[
\int x^3\ln(x)\,dx
\]
\[
u = \ln(x) \quad \Rightarrow \quad du = \frac{1}{x}\,dx
\]
\[
dv = x^3\,dx \quad \Rightarrow \quad v = \frac{x^4}{4}
\]
\[
\int x^3\ln(x)\,dx
= \frac{x^4}{4}\ln(x) - \int \frac{x^4}{4}\cdot\frac{1}{x}\,dx
\]
\[
= \frac{x^4}{4}\ln(x) - \frac{1}{4}\int x^3\,dx
\]
\[
= \frac{x^4}{4}\ln(x) - \frac{x^4}{16} +C
\]

\color{black}

\item Considere la función $f : \mathbb{R} \to \mathbb{R}$ definida por $f(x) = \dfrac{1 - x^2}{1 + x^2}$ y la región $\mathcal{R}$ definida como el sector del plano comprendido entre $f$ y el eje horizontal $OX$, entre los ceros de $f$.
\begin{enumerate}
    \item Demuestre que existen $A, B \in \mathbb{R}$ tales que, para todo $x \in \mathbb{R}$, $f(x) = A + \dfrac{B}{1 + x^2}$

    \color{blue}
    \textbf{Sol:} Sea $\displaystyle A+\frac{B}{1+x^2}$
\[
A+\frac{B}{1+x^2}=\frac{A(1+x^2)+B}{1+x^2}.
\]
Igualamos a $f(x)$
\[
\dfrac{1 - x^2}{1 + x^2}=\frac{A(1+x^2)+B}{1+x^2} \implies
A(1+x^2)+B=1-x^2.
\]
igualamos coeficientes
\[
\begin{cases}
A= -1,\\
A+B=1,
\end{cases}
\Rightarrow\ B=2.
\]
en efecto
\[
f(x)=-1+\frac{2}{1+x^2}, \qquad \forall x\in\mathbb{R}.
\]
    \color{black}
    
    \item Calcule el área de la región $\mathcal{R}$.

\color{blue}
\textbf{Sol:} Los ceros de $f$ se obtienen de
\[
1-x^2=0 \;\Rightarrow\; x=\pm1.
\]
Como $f(x)\ge 0$ en $[-1,1]$, el área es
\[
A=\int_{-1}^{1} f(x)\,dx.
\]
Usando que $\displaystyle f(x)=-1+\dfrac{2}{1+x^2}$,
\[
A=\int_{-1}^{1}\left(-1+\frac{2}{1+x^2}\right)dx.
\]
\[
A=\left[-x+2\arctan x\right]_{-1}^{1}
\]
\[
A=\big(-1+2\arctan 1\big)-\big(1-2\arctan 1\big)
= -2+4\arctan 1.
\]
como $\arctan 1=\dfrac{\pi}{4}$,
\[
A=\pi-2
\]
\color{black}
    
\item Calcule el volumen del sólido de revolución obtenido al rotar la región $\mathcal{R}$.

\textbf{Hint:} Puede usar que $\displaystyle \int \frac{1}{(1+x^2)^2} \, dx = \frac{1}{2} \left( \frac{x}{x^2+1} + \arctan(x) \right) + C$

\color{blue}
\textbf{Sol:} al tratarse de un sólido de revolución:
\begin{align*}
V(\mathcal{R}) &= \pi \int_{-1}^{1} f^2(x) \, dx = \pi \int_{-1}^{1} \left( \frac{2}{1+x^2} - 1 \right)^2 dx \\
&= \pi \int_{-1}^{1} \left( \frac{4}{(1+x^2)^2} - \frac{4}{1+x^2} + 1 \right) dx \\
&= \pi \left( \frac{4}{2} \left( \frac{x}{x^2+1} + \arctan(x) \right) - 4 \arctan(x) + x \right) \bigg|_{-1}^{1} \\
&= \pi \left( 2 \left( \frac{1}{2} + \frac{\pi}{4} \right) - 4 \frac{\pi}{4} + 1 \right. \\
&\quad - \left. \left( 2 \left( -\frac{1}{2} - \frac{\pi}{4} \right) + 4 \frac{\pi}{4} - 1 \right) \right) \\
&= \pi \left( 1 + \pi/2 - \pi + 1 + 1 + \pi/2 - \pi + 1 \right) \\
&= \pi (4 - \pi) > 0 \quad (\text{¡ojo con eso!})
\end{align*}
\color{black}

\end{enumerate}

    
    \item Sea $f(x) = x\sqrt{1 - x^2}$. Si
    \[
    R = \{(x, y) \in \mathbb{R}^2 \mid 0 \le x \le 1, 0 \le y \le f(x)\}.
    \]
    \begin{enumerate}
        \item Encuentre el área de la región $R$.

\color{blue}
\textbf{Sol:} El área de $R$ es
\[
A=\int_{0}^{1} x\sqrt{1-x^2}\,dx.
\]
cambio de variable $u=1-x^2$, entonces $du=-2x\,dx$, (recordar ajustar los límites de integración).
\[
A=\frac12\int_{0}^{1} \sqrt{u}\,du.
\]
\[
A=\frac12\left[\frac{2}{3}u^{3/2}\right]_{0}^{1}=1/3
\]
\color{black}
        
        \item Encuentre el volumen del sólido de revolución que se obtiene al rotar la región $R$ en torno al eje $OX$.

\color{blue}
\textbf{Sol:} Al rotar $R$ alrededor del eje $OX$ usamos
\[
V=\pi\int_{0}^{1} \bigl(f(x)\bigr)^2\,dx.
\]
con $f(x)=x\sqrt{1-x^2}$,
\[
\bigl(f(x)\bigr)^2=x^2(1-x^2).
\]
\[
V=\pi\int_{0}^{1}\left(x^2-x^4\right)dx.
\]
\[
V=\pi\left[\frac{x^3}{3}-\frac{x^5}{5}\right]_{0}^{1}.
\]
\[
V=\pi\left(\frac{1}{3}-\frac{1}{5}\right)=\frac{2\pi}{15}.
\]

\color{black}
          
    \end{enumerate}
    \item 
\begin{enumerate}
\item Sean $\displaystyle I = \int \cos(\ln(x))dx$ y $\displaystyle J = \int \sen(\ln(x))dx$. Usando integración por partes, plantee un sistema lineal que permita calcular $I$ y $J$. Calcule $I$ y $J$.

\color{blue}
\textbf{Sol:} Definimos
\[
I=\int \cos(\ln x)\,dx, \qquad J=\int \sen(\ln x)\,dx.
\]
Para $I$, integramos por partes con
\[
u=\cos(\ln x), \quad dv=dx,
\]
\[
du=-\frac{\sen(\ln x)}{x}dx, \quad v=x.
\]
\[
I=x\cos(\ln x)+\int \sen(\ln x)\,dx=x\cos(\ln x)+J.
\]
Para $J$, integramos por partes con
\[
u=\sen(\ln x), \quad dv=dx,
\]
\[
du=\frac{\cos(\ln x)}{x}dx, \quad v=x.
\]
\[
J=x\sen(\ln x)-\int \cos(\ln x)\,dx=x\sen(\ln x)-I.
\]
obtenemos el sistema lineal:
\[
\begin{cases}
I - J = x\cos(\ln x),\\
I + J = x\sen(\ln x).
\end{cases}
\]
resolvemos
\[
I=\frac{x}{2}\big(\sen(\ln x)+\cos(\ln x)\big)+C,
\]
\[
J=\frac{x}{2}\big(\sen(\ln x)-\cos(\ln x)\big)+C.
\]
\color{black}

\item Deducir una fórmula de recurrencia para $\displaystyle I_{m,n} = \int x^m (\ln(x))^n dx$. Use la fórmula para calcular $\displaystyle \int x^2 \ln x$.

\color{blue}
\textbf{Sol:} Sea
\[
I_{m,n}=\int x^m(\ln x)^n\,dx,\qquad m\neq -1.
\]
Integramos por partes tomando
\[
u=(\ln x)^n,\quad dv=x^m dx,
\]
de donde
\[
du=\frac{n(\ln x)^{\,n-1}}{x}dx,\quad v=\frac{x^{m+1}}{m+1}.
\]
\[
I_{m,n}=\frac{x^{m+1}}{m+1}(\ln x)^n-\frac{n}{m+1}\int x^m(\ln x)^{n-1}dx.
\]
Obtenemos la fórmula de recurrencia:
\[
\,I_{m,n}=\frac{x^{m+1}}{m+1}(\ln x)^n-\frac{n}{m+1}I_{m,n-1}\, .
\]
para $\displaystyle \int x^2\ln x\,dx$, tenemos $m=2$, $n=1$.
\[
I_{2,1}=\frac{x^3}{3}\ln x-\frac{1}{3}I_{2,0}.
\]
Como
\[
I_{2,0}=\int x^2dx=\frac{x^3}{3},
\]
\[
\int x^2\ln x\,dx=\frac{x^3}{3}\ln x-\frac{x^3}{9}+C.
\]
\color{black}

\end{enumerate}
\newpage
\item Dadas las curvas $y = mx$ y $y = x^2$, considere la región limitada por ambas curvas y encuentre el valor de $m > 0$, para que los volúmenes de los sólidos obtenidos al rotar la región definida en torno al eje $OX$ y al eje $OY$, sean iguales. Luego, calcule el área encerrada entre ambas curvas.

\color{blue}
\textbf{Sol:} \begin{figure}[H]
    \centering
    \includegraphics[width=0.5\linewidth]{Imagenes/grafica_integral.png}
\end{figure}
Las curvas $y=mx$ y $y=x^2$ se intersectan cuando
\[
x^2=mx \;\Rightarrow\; x=0,\; x=m.
\]
Volumen al rotar en torno a $OX$ (método de anillos):
\[
V_{OX}=\pi\int_{0}^{m}\big[(mx)^2-(x^2)^2\big]dx
=\pi\int_{0}^{m}(m^2x^2-x^4)dx.
\]
\[
V_{OX}=\pi\left[\frac{m^2x^3}{3}-\frac{x^5}{5}\right]_{0}^{m}
=\pi m^5\left(\frac{1}{3}-\frac{1}{5}\right)
=\frac{2\pi m^5}{15}.
\]
Volumen al rotar en torno a $OY$ (método de capas cilíndricas):
\[
V_{OY}=2\pi\int_{0}^{m}x(mx-x^2)\,dx
=2\pi\int_{0}^{m}(mx^2-x^3)dx.
\]
\[
V_{OY}=2\pi\left[\frac{mx^3}{3}-\frac{x^4}{4}\right]_{0}^{m}
=2\pi m^4\left(\frac{1}{3}-\frac{1}{4}\right)
=\frac{\pi m^4}{6}.
\]
Igualando volúmenes $V_{OX}=V_{OY}$:
\[
\frac{2\pi m^5}{15}=\frac{\pi m^4}{6}
\;\Rightarrow\;
m=\frac{5}{4}.
\]
Área encerrada entre las curvas:
\[
A=\int_{0}^{m}(mx-x^2)\,dx.
\]
Con $m=\tfrac{5}{4}$,
\[
A=\left[\frac{mx^2}{2}-\frac{x^3}{3}\right]_{0}^{5/4}
=\frac{125}{384}.
\]



\color{black}



    \item Dada la elipse de ecuación
    \[
    \frac{x^2}{2} + y^2 = 1,
    \]
    encuentre el volumen del manto generado al rotar esta elipse en torno al eje $OX$ entre $x = -1$ y $x = 1$.

\color{blue}
\textbf{Sol:} 
De la elipse
\[
\frac{x^2}{2}+y^2=1
\]
se obtiene
\[
y^2=1-\frac{x^2}{2}.
\]
Al rotar alrededor del eje $OX$ se obtiene
\[
V=\pi\int_{-1}^{1} y^2\,dx
=\pi\int_{-1}^{1}\left(1-\frac{x^2}{2}\right)dx.
\]
\[
V=\pi\left[x-\frac{x^3}{6}\right]_{-1}^{1}.
\]
\[
V=\pi\left[\left(1-\frac{1}{6}\right)-\left(-1+\frac{1}{6}\right)\right]
=\frac{5\pi}{3}
\]
\color{black}

\begin{enumerate}


\item Calcule las siguientes primitivas usando integraci\'on por partes
\begin{enumerate}
    \item $\displaystyle \int x \cos(2x) dx$

\color{blue}
\textbf{Sol:} $$
\begin{aligned}
\text{Sea } u &= x \quad \Rightarrow \quad du = dx \\
dv &= \cos(2x)\,dx \quad \Rightarrow \quad v = \frac{1}{2} \sen(2x) \\[6pt]
\int x \cos(2x)\,dx 
&= uv - \int v\,du \\
&= x \cdot \frac{1}{2} \sen(2x) - \int \frac{1}{2} \sen(2x)\,dx \\[6pt]
&= \frac{x}{2} \sen(2x) - \frac{1}{2} \int \sen(2x)\,dx \\
&= \frac{x}{2} \sen(2x) + \frac{1}{4} \cos(2x) + C
\end{aligned}
$$
\color{black}

    
    \item $\displaystyle \int x \arctan(x) dx$  

\color{blue}
\textbf{Sol:} $$
\begin{aligned}
\text{Sea } u &= \arctan(x) \ \Rightarrow \ du = \frac{1}{1 + x^2} \, dx \\
dv &= x\,dx \ \Rightarrow \ v = \frac{x^2}{2} \\[6pt]
\int x \arctan(x)\,dx 
&= uv - \int v\,du \\
&= \frac{x^2}{2} \cdot \arctan(x) - \int \frac{x^2}{2} \cdot \frac{1}{1 + x^2} \, dx \\[6pt]
&= \frac{x^2}{2} \arctan(x) - \frac{1}{2} \int \frac{x^2}{1 + x^2} \, dx \\
&= \frac{x^2}{2} \arctan(x) - \frac{1}{2} \int \left(1 - \frac{1}{1 + x^2}\right) dx \\
&= \frac{x^2}{2} \arctan(x) - \frac{1}{2} \left( x - \arctan(x) \right) + C \\
&= \frac{x^2}{2} \arctan(x) - \frac{x}{2} + \frac{1}{2} \arctan(x) + C \\[6pt]
&= \left( \frac{x^2}{2} + \frac{1}{2} \right) \arctan(x) - \frac{x}{2} + C
\end{aligned}
$$
\color{black}

    
    \item $\displaystyle \int \ln^2(x) dx$ 

\color{blue}
\textbf{Sol:} $$
\begin{aligned}
\text{Sea } u &= \ln^2(x) \ \Rightarrow \ du = \frac{2\ln(x)}{x}\,dx \\
dv &= dx \ \Rightarrow \ v = x \\[6pt]
\int \ln^2(x)\,dx 
&= uv - \int v\,du \\
&= x \ln^2(x) - \int x \cdot \frac{2\ln(x)}{x}\,dx \\
&= x \ln^2(x) - 2 \int \ln(x)\,dx \\
&= x \ln^2(x) - 2(x \ln(x) - x) + C \\
&= x \ln^2(x) - 2x \ln(x) + 2x + C
\end{aligned}
$$

\color{black}

    
    \item $\displaystyle \int x^3 \ln(x) dx$

\color{blue}
\textbf{Sol:} $$
\begin{aligned}
\text{Sea } u &= \ln(x) \ \Rightarrow \ du = \frac{1}{x} dx \\
dv &= x^3\,dx \ \Rightarrow \ v = \frac{x^4}{4} \\[6pt]
\int x^3 \ln(x)\,dx 
&= uv - \int v\,du \\
&= \frac{x^4}{4} \ln(x) - \int \frac{x^4}{4} \cdot \frac{1}{x}\,dx \\
&= \frac{x^4}{4} \ln(x) - \frac{1}{4} \int x^3\,dx \\
&= \frac{x^4}{4} \ln(x) - \frac{1}{4} \cdot \frac{x^4}{4} + C \\
&= \frac{x^4}{4} \ln(x) - \frac{x^4}{16} + C
\end{aligned}
$$
\color{black}

    
    \item $\displaystyle \int \sen(x) e^{2x} dx$

\color{blue}
\textbf{Sol:}  
\begin{align*}
\text{Sea } & u = e^{2x}, \quad dv = \sen(x) \, dx \\
& du = 2 e^{2x} \, dx, \quad v = -\cos(x) \\
\end{align*}

\vspace{-1cm}
Reemplazando
\vspace{-0.5cm}

\begin{align*}
\int \sen(x) e^{2x} \, dx &= - e^{2x} \cos(x) + 2 \underset{I}{\int e^{2x} \sen(x) \, dx} \\
\end{align*}
\vspace{-1cm}
Resolviendo $I$
\vspace{-0.5cm}
\begin{align*}
\text{Sea} \ & u = e^{2x}, \quad dv = \cos(x) \, dx \\
& du = 2 e^{2x} \, dx, \quad v = \sen(x) \\
 I &= \sen(x) e^{2x} - 2 \int e^{2x} \sen(x) \, dx \\
\therefore \quad \int \sen(x) e^{2x} \, dx &= - e^{2x} \cos(x) + 2 \left( \sen(x) e^{2x} -2 \int e^{2x} \sen(x) \, dx\right) \\
\int \sen(x) e^{2x} \, dx &= - e^{2x} \cos(x) + 2 \sen(x) e^{2x} -4 \int e^{2x} \sen(x) \, dx \\
5\int \sen(x) e^{2x} \, dx &= - e^{2x} \cos(x) + 2 \sen(x) e^{2x} \\
\text{Finalmente} \quad \int \sen(x) e^{2x} \, dx &= - \frac{e^{2x} \cos(x)}{5} + \frac{2}{5} \sen(x) e^{2x} + C
\end{align*}
\color{black}


    
    \item $\displaystyle \int \frac{xe^{x}}{(1+x)^{2}}dx. $

\color{blue}
\textbf{Sol:} $$
\begin{aligned}
\text{Sea } t &= 1 + x \quad \Rightarrow \quad dt = dx, \quad x = t - 1 \\[6pt]
\int \frac{x e^{x}}{(1+x)^2} dx 
&= \int \frac{(t - 1) e^{t - 1}}{t^2} dt \\
&= \int \frac{(t - 1)}{t^2} e^{t - 1} dt = \int \left( \frac{t}{t^2} - \frac{1}{t^2} \right) e^{t - 1} dt \\
&= \int \left( \frac{1}{t} - \frac{1}{t^2} \right) e^{t - 1} dt = e^{-1} \int \left( \frac{1}{t} - \frac{1}{t^2} \right) e^{t} dt \\
&= e^{-1} \int \left( \frac{1}{t} e^{t} - \frac{1}{t^2} e^{t} \right) dt
\end{aligned}
$$
\color{black}

    
\end{enumerate}
\end{enumerate}


\item Seguramente recuerda la fórmula para el volumen
de una pirámide. Entonces demuestrela: 
\begin{figure}[H]
    \centering
    \includegraphics[width=0.8\linewidth]{Imagenes/piramides.png}
    \caption{pirámide de base cuadrada (a. izq) y pirámide de base irregular (b. der)}
\end{figure}

\color{blue}
\textbf{Sol:} ver el siguiente video https://www.youtube.com/watch?v=6rynOsZppEE (si hago la pauta de esto no se va a entender).
\color{black}

\item Considere la función $f : [0, \frac{1}{2}] \to \mathbb{R}$ definida por
$$
f(x) = -\ln(1 - x^2)
$$
Sea $A$ la región del plano limitada por la curva $f(x)$ y el eje $OX$ para $x \in [0, \frac{1}{2}]$.
\begin{enumerate}
    \item[(a)] Calcule el área de la región $A$.

\color{blue}
    \textbf{Sol:} El área de la región $A$ viene dada por
\[
A=\int_{0}^{1/2} -\ln(1-x^2)\,dx.
\]
Aplicamos integración por partes
\[
u=-\ln(1-x^2), \qquad dv=dx,
\implies
du=\frac{2x}{1-x^2}\,dx, \qquad v=x.
\]
\[
\int -\ln(1-x^2)\,dx
= -x\ln(1-x^2)-\int \frac{2x^2}{1-x^2}\,dx.
\]
tenemos que
\[
\frac{2x^2}{1-x^2}=-2+\frac{2}{1-x^2},
\]
\[
\int -\ln(1-x^2)\,dx
= -x\ln(1-x^2)+2x-2\,\operatorname{artanh}(x)+C.
\]

Evaluando entre $0$ y $\tfrac12$:
\[
A=\left[-x\ln(1-x^2)+2x-2\,\operatorname{artanh}(x)\right]_{0}^{1/2}.
\]
\[
\operatorname{artanh}\!\left(\tfrac12\right)=\tfrac12\ln 3,
\]
\[
A=-\tfrac12\ln\!\left(\tfrac34\right)+1-\ln 3
=1+\ln 2-\tfrac32\ln 3.
\]
Por lo tanto, el área de la región $A$ es
\[
A=1+\ln 2-\tfrac{3}{2}\ln 3.
\]
    \color{black}
    \item[(b)] Calcule el volumen del sólido de revolución que se obtiene al rotar la región $A$ en torno al eje $OY$.

\color{blue}
\textbf{Sol:} \[
\begin{aligned}
V_{OY} &= 2\pi \int_{a}^{b} x f(x) \, dx \implies V_{OY} = 2\pi \int_{0}^{1/2} -x \ln(1-x^2) \, dx \\
&\quad \text{c.v. } u = 1-x^2 \\
&\quad du = -2x \, dx \implies V_{OY} = \pi \int_{1}^{3/4} \ln(u) \, du = \dots \\
&\quad u(0)=1, \ u(1/2)=3/4 \\
\dots &= -\pi \Big( u \ln u - u \Big) \bigg|_{1}^{3/4} \\
&= \pi \Big( u - u \ln(u) \Big) \bigg|_{3/4}^{1} \\
&= \pi \left( 1 - \underbrace{1 \ln(1)}_{0} - \frac{3}{4} + \frac{3}{4} \ln\left(\frac{3}{4}\right) \right) \\
&= \pi \left( \frac{1}{4} + \frac{3}{4} \ln\left(\frac{3}{4}\right) \right) > 0
\end{aligned}
\]
\color{black}

    
\end{enumerate}
\item Calcule el \'area encerrada entre las curvas $y = 15 - x^2$ e $y = x^4 - 5$.

\color{blue}
\textbf{Sol:} hallamos los puntos de intersección
\[
15 - x^2 = x^4 - 5,
\implies
x^4 + x^2 - 20 = 0.
\]
Sea $u = x^2$
\[
u^2 + u - 20 = 0,
\]
con $u = 4$ y $u = -5$. Luego como $u = x^2 \ge 0$
\[
x^2 = 4 \Rightarrow x = \pm 2.
\]
luego en el intervalo $[-2,2]$ la parábola $y=15-x^2$ está por encima de
$y=x^4-5$, esto pues considerando la diferencia entre las funciones
\[
(15-x^2)-(x^4-5)=20-x^2-x^4.
\implies \text{Sea }
g(x)=20-x^2-x^4.
\]
En $[-2,2]$ se tiene
\[
g(x)\ge 0,
\]
se puede comprobar con la tabla de signosla inferior es. Como la función inferior es
$y = x^4 - 5$. Luego el área buscada es
\[
A = \int_{-2}^{2} \left[(15 - x^2) - (x^4 - 5)\right]\,dx
= \int_{-2}^{2} (20 - x^2 - x^4)\,dx.
\]
Como el integrando es una función par,
\[
A = 2 \int_{0}^{2} (20 - x^2 - x^4)\,dx.
\implies 
\int (20 - x^2 - x^4)\,dx
= 20x - \frac{x^3}{3} - \frac{x^5}{5}.
\]
\[
A = 2\left[20(2) - \frac{2^3}{3} - \frac{2^5}{5}\right]
= 2\left(40 - \frac{8}{3} - \frac{32}{5}\right).
\]
\[
A = \frac{928}{15}.
\]
\color{black}



\item Determine el \'area entre las curvas $y = f(x) = x^3 - x^2$ e $y = g(x) = 2x$ en el intervalo $[-2, 2]$. %con pauta examen 2023-2

\color{blue}
\textbf{Sol:} Primero buscamos los puntos donde las curvas se intersectan resolviendo la ecuación
\[
f(x)=x^3-x^2=2x=g(x),
\]
la cual se factoriza como
\[
f(x)-g(x)=x(x-2)(x+1)=0.
\]
Así, las curvas se intersectan en $x=-1$, $x=0$ y $x=2$. La función $g$ es mayor que $f$ en los intervalos $[-2,-1]$ y $[0,2]$, por lo que el área buscada es
\[
A=\int_{-2}^{2} |f(x)-g(x)|\,dx
= \int_{-2}^{-1} |f(x)-g(x)|\,dx
+ \int_{-1}^{0} |f(x)-g(x)|\,dx
+ \int_{0}^{2} |f(x)-g(x)|\,dx.
\]
Como $g(x)\ge f(x)$ en los intervalos considerados, se tiene
\[
A=\int_{-2}^{-1} (g(x)-f(x))\,dx
- \int_{-1}^{0} (g(x)-f(x))\,dx
+ \int_{0}^{2} (g(x)-f(x))\,dx.
\]
Una primitiva particular de $g(x)-f(x)$ es
\[
F(x)=-\frac{x^4}{4}+\frac{x^3}{3}+x^2.
\]
Por el Teorema Fundamental del Cálculo,
\[
A=\left[F(x)\right]_{-2}^{-1}
-\left[F(x)\right]_{-1}^{0}
+\left[F(x)\right]_{0}^{2}.
\]
\[
A=F(-1)-F(-2)-F(0)+F(-1)+F(2)-F(0)
=F(2)-F(-2)+2F(-1)-2F(0).
\]
Finalmente, usando
\[
F(-2)=-\frac{8}{3}, \quad
F(2)=\frac{8}{3}, \quad
F(0)=0, \quad
F(-1)=\frac{5}{12},
\]
se obtiene
\[
A=\frac{16}{3}+\frac{5}{6}=\frac{37}{6}.
\]

\color{black}

\item Considere las funciones
$$ f(x) = x^{2}-1 \quad \text{y} \quad g(x) = -x^{2}+1. $$
Se pide:
\begin{enumerate}
\item Bosquejar las gr\'aficas de $f$ y $g$ y determinar los puntos en donde se intersectan.

\color{blue}
\begin{figure}[H]
    \centering    \includegraphics[width=0.35\linewidth]{Imagenes/bosquejo.png}
\end{figure}
Donde la curva verde es la gráfica de $f$ y la azul es la de $g$. Para hallar
los puntos de intersección se iguala $f(x)=g(x)$, de donde
\[
x^2-1=-x^2+1 \;\Rightarrow\; x^2=1 \;\Rightarrow\; x=\pm 1.
\]
Como $f(\pm 1)=g(\pm 1)=0$, las gráficas de $f$ y $g$ se intersectan en los
puntos
\[
(-1,0) \quad \text{y} \quad (1,0).
\]

\color{black}

\item Calcular el área encerrada (delimitada entre los puntos de intersección) por las funciones $f(x)$ y $g(x)$.

\color{blue}
\textbf{Sol:} El área es
\[
\text{Área}
= \int_{-1}^{1} (g(x)-f(x))\,dx
= \int_{-1}^{1} \bigl(-x^2+1-(x^2-1)\bigr)\,dx.
\]
\[
\text{Área}
= \int_{-1}^{1} (-2x^2+2)\,dx.
\]
\[
\text{Área}
= \left(-\frac{2}{3}x^3+2x\right)\Big|_{-1}^{1}
= \left(-\frac{2}{3}+2\right)-\left(\frac{2}{3}-2\right)
= \frac{8}{3}.
\]

\color{black}
\end{enumerate}

\begin{enumerate}
\item 
Cálcule
\begin{enumerate}
    \item El \'area que encierra la funci\'on $\cos(t)$ en $0$ y $2\pi$.

\color{blue}
\textbf{Sol:} La función \(\cos(t)\) es positiva en \([0, \frac{\pi}{2}), (\frac{3\pi}{2},\frac{5\pi}{2})\) y negativa en \((\frac{\pi}{ 2},\frac{3\pi}{2}).(\frac{5\pi}{2},2\pi)\), tenemos que tener cuidado de sumar área negativa, por lo tanto, como en dichos intervalos, el área es simétrica con respecto al eje $x$, tenemos que 
el área es
\[
A = 4 \int_0^{\pi/2} \cos(t) dt = 2 \left( \sen\left(\frac{\pi}{2}\right) - \sen(0) \right) = 4 \cdot 1 = 4.
\]
\begin{figure}[H]
    \centering
    \includegraphics[width=0.7\linewidth]{Imagenes/integral_cos.png}
    \caption{Integral pedida, tenemos que considerar $\displaystyle 4 \int_0^{\pi/2} \cos(t) dt$ para que no se nos cancele el área.}
    \label{fig:placeholder}
\end{figure}
\color{black}

    \item El \'area de una circunferencia de radio $R$.

\color{blue}
\textbf{Sol:} 
Tenemos que queremos calcular el área del círculo de radio $R$ centrado en el origen:
\[
D = \{(x,y) \in \mathbb{R}^2 : x^2 + y^2 \leq R^2 \}.
\]
Para esto consideramos su ecuación de circunferencia en $\R^2$.
$$
x^2 + y^2 = R^2 \Rightarrow
y = \sqrt{R^2 - x^2}
$$
En este último despeje, tenemos la mitad de la circuferencia a la que queremos calcularle el área. Por lo tanto el área total va a ser
$$A = 2 \int_{-R}^{R} \sqrt{R^2 - x^2} \, dx
$$
Esta integral es del tipo trigonométrica. Luego
\begin{align*}
\text{Sea } x &= R \sen(\theta) \\
dx &= R \cos(\theta) \, d\theta
\end{align*}
Cambiamos los límites de Integración:
\begin{multicols}{2}
\begin{itemize}
    \item Límite Superior ($x=R$):
    \begin{align*}
    x = R &\Rightarrow R = R \sen(\theta) \\
          &\Rightarrow 1 = \sen(\theta) \\
          &\Rightarrow \theta = \frac{\pi}{2}
    \end{align*}
    \item Límite Inferior ($x=-R$):
    \begin{align*}
    x = -R &\Rightarrow -R = R \sen(\theta) \\
           &\Rightarrow -1 = \sen(\theta) \\
           &\Rightarrow \theta = -\frac{\pi}{2}
    \end{align*}
\end{itemize}
\end{multicols}
\begin{multicols}{2}
\begin{align*}
\therefore \, A &= 2 \int_{-\pi/2}^{\pi/2} \sqrt{R^2 - R^2 \sen^2(\theta)} \cdot R \cos(\theta) \, d\theta \\
&= 2 \int_{-\pi/2}^{\pi/2} \sqrt{R^2(1 - \sen^2(\theta))} \cdot R \cos(\theta) \, d\theta \\
&= 2 \int_{-\pi/2}^{\pi/2} \sqrt{R^2 \cos^2(\theta)} \cdot R \cos(\theta) \, d\theta \\
&= 2 \int_{-\pi/2}^{\pi/2} R \cos(\theta) \cdot R \cos(\theta) \, d\theta \\
&= 2 \int_{-\pi/2}^{\pi/2} R^2 \cos^2(\theta) \, d\theta\\
&= 2 R^2 \int_{-\pi/2}^{\pi/2} \cos^2(\theta) \, d\theta \\
&= 2 R^2 \int_{-\pi/2}^{\pi/2} \frac{1 + \cos(2\theta)}{2} \, d\theta \\
&= 2 R^2 \left[ \frac{1}{2}\theta + \frac{1}{4}\sen(2\theta) \right]_{-\pi/2}^{\pi/2} \\
&= 2 R^2 \left[ \left(\frac{\pi}{4} + \frac{1}{4}\sen(\pi)\right) - \left(-\frac{\pi}{4} + \frac{1}{4}\sen(-\pi)\right) \right] \\
&= 2 R^2 \left[ \left(\frac{\pi}{4} + 0\right) - \left(-\frac{\pi}{4} + 0\right) \right] \\
&= 2 R^2 \left(\frac{\pi}{4} - \left(-\frac{\pi}{4}\right)\right) \\
&= 2 R^2 \left(\frac{\pi}{2}\right) = \pi R^2
\end{align*}
\includegraphics[width=1\linewidth]{Imagenes/c_variable_y_circulo.png}
\end{multicols}
\color{black}

    \item El volumen de una esfera de radio $R$.

\color{blue}
\textbf{Sol:} Consideramos la semi circunferencia por encima del eje $x$ de la esfera dada por
\[
y = \sqrt{R^2 - x^2}, \quad x \in [-R, R].
\]
\begin{figure}[H]
    \centering
    \includegraphics[width=0.5\linewidth]{Imagenes/semi_circunferencia.png}
    \caption{Semi circunferencia en $\R^2$ para $R=1$}
    \label{fig:placeholder}
\end{figure}
Al girar esta curva alrededor del eje \(x\), el volumen del sólido de revolución es
\[
V = \pi \int_{-R}^R \left( \sqrt{R^2 - x^2} \right)^2 dx = \pi \int_{-R}^R (R^2 - x^2) \, dx.
\]
\[
V = \pi \left[ R^2 x - \frac{x^3}{3} \right]_{-R}^R = \pi \left( R^3 - \frac{R^3}{3} - (-R^3 + \frac{R^3}{3}) \right) = \pi \left( \frac{2R^3}{3} + \frac{2R^3}{3} \right) = \frac{4\pi R^3}{3}.
\]
\begin{multicols}{2}
\begin{figure}[H]
    \centering
    \includegraphics[width=0.7\linewidth]{Imagenes/semi_circunferencia_2.png}
    \label{fig:placeholder}
\end{figure}
\begin{figure}[H]
    \centering
    \includegraphics[width=0.8\linewidth]{Imagenes/semi_circunferencia_3.png}
    \label{fig:placeholder}
\end{figure}
\end{multicols}
En estas imagenes vemos que la idea del solido de revolución es partir con la media circunferencia, rotarla en el eje $x$ y proyectarla en el eje $z$ para así obtener el volumen de la esfera en $\R^3$, que para este ejemplo consideré el radio $R=1$. 
\color{black}

    \item El volumen de una pir\'amide de base cuadrada de lado $a$ y de altura $h$.

\color{blue}
\textbf{Sol:} Propuesto, ver el video de matefacil en puesto en la descripción del pdf.
\color{black}
    
    \item Calcular el \'area del toro, es decir del cuerpo generado al rotar una circunferencia de radio $r$ y de centro en $(0,R)$ ($R>r$), en torno al eje $OX$ (como en la figura).
    
\end{enumerate}
\begin{figure}[H]
    \centering
    \includegraphics[width=0.5\linewidth]{Imagenes/untoro.jpg}
    \caption{Un toro}
    \label{fig:placeholder}
\end{figure}
\color{blue}
\textbf{Sol:} Sea una circunferencia de radio \( r \) con centro en el punto \( (0, R) \), donde \( R > r \).

La ecuación de la circunferencia es:
\[
x^2 + (y - R)^2 = r^2.
\]

Despejamos \( y \):
\[
y(x) = R \pm \sqrt{r^2 - x^2}, \quad x \in [-r, r].
\]

Para obtener el área superficial generada al rotar esta curva alrededor del eje \(OX\), usamos la fórmula del área de superficie de revolución:
\[
A = 2\pi \int_{-r}^r y(x) \sqrt{1 + \left( \frac{dy}{dx} \right)^2} \, dx.
\]

Como la curva es simétrica respecto al eje \( y = R \), basta con usar solo la mitad superior:
\[
y(x) = R + \sqrt{r^2 - x^2}, \quad \text{entonces } \frac{dy}{dx} = \frac{-x}{\sqrt{r^2 - x^2}}.
\]

Entonces:
\[
\left( \frac{dy}{dx} \right)^2 = \frac{x^2}{r^2 - x^2}, \quad
1 + \left( \frac{dy}{dx} \right)^2 = \frac{r^2}{r^2 - x^2}.
\]

Sustituimos en la fórmula del área:
\[
A = 2\pi \int_{-r}^r \left(R + \sqrt{r^2 - x^2}\right) \cdot \sqrt{\frac{r^2}{r^2 - x^2}} \, dx.
\]

\[
A = 2\pi r \int_{-r}^r \left( \frac{R + \sqrt{r^2 - x^2}}{\sqrt{r^2 - x^2}} \right) dx = 2\pi r \int_{-r}^r \left( \frac{R}{\sqrt{r^2 - x^2}} + 1 \right) dx.
\]
Ahora resolvemos cada parte:
\begin{multicols}{2}
\begin{itemize}
    \item \(\displaystyle \int_{-r}^r \frac{dx}{\sqrt{r^2 - x^2}} = \pi\)
    \item \(\displaystyle \int_{-r}^r dx = 2r\)
\end{itemize}
\end{multicols}
Entonces:
\[
A = 2\pi r \left( R \cdot \pi + 2r \right) = 2\pi r (R\pi + 2r).
\]
\color{black}


    
\end{enumerate}
\end{enumerate}
\end{document}